\section{Collect, lessons, Gospel, Creed (nn.~1026--1029)}

Between the \latincue{Gloria} and the offertory the missal prints almost no
text, because almost everything here is proper. What it prints instead is a
sequence of cues and two short prayers, and the cues are worth reading closely,
because they tell a reader what the Roman rite thinks a lesson is.

\subsection*{The collect and the lessons}

The \work{Ordo Missae} says only this: the priest kisses the altar, turns and
says \latincue{Dominus vobiscum}, then \latincue{Oremus} and the orations, one
or more as the Office requires; then follow the Epistle, gradual, tract or
Alleluia with its verse, or sequence, as the season or the character of the
Mass demands.

Three features of that sentence are rubrically loaded. The plural
\latincue{orationes, unam aut plures} points to the 1960 code's reduction of
the old accumulation of collects: nn.~433--435 fix how many orations may be
said on each class of day and forbid exceeding three under any pretext. The
phrase \latincue{ut Ordo Officii postulat} makes the Mass depend on the
calendar rather than on the celebrant's choice. And the list of chants ---
gradual, tract, Alleluia, sequence --- is given as alternatives keyed to
season, so that in this rite the chant between the lessons is not one item with
seasonal decorations but four different items.

\begin{rubricnote}{The lessons in the general rubrics, nn.~466--474}
After the orations the Epistle is said, and \latincue{Deo gratias} is answered.
One lesson precedes the Epistle on the Ember Wednesdays, on the Wednesday of
the fourth week of Lent and on Wednesday in Holy Week; five precede it on the
Ember Saturdays, with \latincue{Deo gratias} after each except the lesson from
Daniel. At the Gospel, \latincue{Dominus vobiscum} is said, then
\latincue{Sequentia} or \latincue{Initium sancti Evangelii secundum N.}, to
which is answered \latincue{Gloria tibi, Domine}, and at the end \latincue{Laus
tibi, Christe}. In Holy Week, before the Passion, none of that is said: the
announcement is simply \latincue{Passio Domini nostri Iesu Christi secundum
N.}, and there is no \latincue{Laus tibi, Christe} at the end. In sung Masses
the celebrant omits whatever the deacon, subdeacon or lector sings by right of
their own office. After the Gospel, especially on Sundays and holydays of
obligation, there is to be a short homily; and if the homily is given by a
priest other than the celebrant, it must not be superimposed on the celebration
--- the Mass is suspended and resumed only when the homily is finished.
\end{rubricnote}

The last of these is the most striking rubric in the 1960 code, and it is easy
to pass over. It refuses to let preaching happen \emph{during} Mass when
another man is preaching: rather than have two liturgical actions running at
once, the Mass stops. Whatever one makes of the practice, the rubric is
evidence that the compilers of this code thought of the Mass as a single
action with a single principal agent, and were willing to interrupt it rather
than blur that.

\subsection*{\latincue{Munda cor meum} (nn.~1026--1028)}

\begin{appointed}{1026--1027 $\cdot$ At solemn Mass}
\textit{Focused source locator.} The appointed unit is identified by this incipit and marginal number; only phrases required by the analysis below are quoted.
\end{appointed}

\begin{appointed}{1028 $\cdot$ At a Mass without sacred ministers}
\textit{Focused source locator.} The appointed unit is identified by this incipit and marginal number; only phrases required by the analysis below are quoted.
\end{appointed}

The prayer is an exegesis of Isaias 6, 6--7 turned into a petition. In the
prophet's vision a seraph takes a live coal from the altar with tongs, touches
his mouth, and says that his iniquity is taken away and his sin cleansed. The
prayer names the episode (\latincue{qui labia Isaiae prophetae calculo mundasti
ignito}), and then asks for the same thing by another means: not by a coal but
\latincue{tua grata miseratione}. The reason the cleansing is asked for is
stated exactly: \latincue{ut sanctum Evangelium tuum digne valeam nuntiare}, in
order to be able to announce, worthily, the Gospel. The verb is
\latincue{nuntiare}, to report or announce, not to explain.

The two forms at nn.~1027 and 1028 differ only in person, and the difference is
the whole of the rite's teaching about who reads the Gospel. At solemn Mass the
deacon asks a blessing from the priest --- \latincue{Iube, domne, benedicere},
with \latincue{domne} addressed to a man --- and the priest answers in the
second person: may the Lord be in \emph{your} heart and on \emph{your} lips. At
a Mass without ministers the priest asks the same blessing of God ---
\latincue{Iube, Domine, benedicere}, with \latincue{Domine} --- and says the
same words of himself in the first person. The single letter that separates
\latincue{domne} from \latincue{Domine} carries the difference between a
delegated office and an unshared one.

After the Gospel the priest kisses the book and says \latincue{Per evangelica
dicta deleantur nostra delicta}. It is not printed as a prayer of the people;
the 1861 hand missal has the congregation join in it as a devotion, which is a
useful reminder that a lay book of that period recorded custom as well as text.

\subsection*{The Creed (n.~1029)}

\begin{appointed}{1029 $\cdot$ \latincue{Credo in unum Deum}}
\textit{Focused source locator.} The appointed unit is identified by this incipit and marginal number; only phrases required by the analysis below are quoted.
\end{appointed}

\begin{rubricnote}{Gesture and occasion}
The rubric printed with the text directs the priest to extend, raise and join
his hands at the opening, to bow his head to the cross at \latincue{Deum} and
again at \latincue{Iesum Christum}, to genuflect at \latincue{Et incarnatus
est} and remain so until \latincue{Et homo factus est} has been said, and to
sign himself from forehead to breast at \latincue{Et vitam venturi saeculi}.
The Creed is printed, like the \latincue{Gloria}, first as chant incipits and
then as text. General rubrics nn.~475--476 fix the occasions: it is said on
every Sunday even when the Office yields to a feast, on first-class feasts and
at first-class votive Masses, on second-class feasts of the Lord and of Our
Lady, through the octaves of Christmas, Easter and Pentecost, and on the feasts
of Apostles and Evangelists together with the Chair of Saint Peter and Saint
Barnabas. It is not said at the chrism Mass or on Maundy Thursday, at the Easter
Vigil, on other second-class feasts, at second-class votive Masses, at festive
and votive Masses of the third and fourth class, on account of a commemoration,
or at Requiems.
\end{rubricnote}

\begin{englishwitness}{xxi--xxii}
``I believe in one God, the Father Almighty, Maker of heaven and earth, and of
all things visible and invisible. And in one Lord Jesus Christ, the only
begotten Son of God, and born of the Father before all ages. God of God; Light
of Light; true God of true God; begotten not made; consubstantial to the
Father, by whom all things were made \dots\ And in the Holy Ghost, the Lord and
Giver of life, who proceedeth from the Father and the Son; who, together with
the Father and the Son, is adored and glorified; who spoke by the prophets. And
one holy Catholic and apostolic Church. I confess one baptism for the remission
of sins. And I expect the resurrection of the dead, and the life of the world to
come. Amen.'' The 1861 book sets \latincue{ET HOMO FACTUS EST} and ``AND WAS
MADE MAN'' in capitals and footnotes them: ``At these words the assistants
kneel down to adore God for the ineffable mystery of the incarnation.''
\end{englishwitness}

This is the Latin liturgical recension of the Nicene-Constantinopolitan Creed.
It receives the conciliar professions of Nicaea and Constantinople and includes
the later Latin \latincue{Filioque}; it is not a text composed in this complete
form by one general council. Its clauses are profession rather than petition,
and several are polemical in origin:
\latincue{Deum verum de Deo vero}, \latincue{genitum, non factum},
\latincue{consubstantialem Patri} are not three ways of saying one thing but
three successive closures against a denial. \latincue{Cuius regni non erit
finis} closes another. The Creed's presence at Mass is therefore an act of a
different kind from the \latincue{Gloria}: the Church is not praising but
professing, and the missal marks the change by having the priest bow to the
cross at the divine names rather than to the altar.

Two ceremonial marks accompany this profession. At
\latincue{Et incarnatus est} the whole assembly kneels, and stays kneeling for
one clause only --- through \latincue{et homo factus est}. At
\latincue{Et vitam venturi saeculi} the priest signs himself with the cross.
The rubrics establish these gestures and their locations; the book does not
state a further symbolic intention for them.

\begin{witnesscard}{Historical note: when the Creed entered the Roman Mass}
Berno of Reichenau reports what he saw at Rome in 1014. The Emperor Henry II
noticed that no creed was sung at the Mass of his coronation on 14 February of
that year, whereas he was used to it in Germany; he was told that the Roman
Church had never been stained by heresy and so had no need to recite it; and
Benedict VIII eventually acceded to the Emperor's wish and ordered it sung
after the Gospel at Rome as well. Fortescue records that most authors accept
this account and date the Creed at Rome from 1014, and he names in a footnote
those who do not --- Probst, who thought Damasus introduced it; Mabillon;
Thalhofer; Mart\`ene, who read Berno to mean that before 1014 it was said and
not sung. The dating is therefore probable rather than secure. Fortescue adds
the observation that the Creed's restriction to Sundays and feasts ``is still a
sign that it is not an essential element'', and the 1962 rubrics preserve
exactly that restriction.\footnote{Fortescue, \work{The Mass}, p.~288, read at
the page image; the dissenting authorities are named in the footnotes on the
same page and are reproduced here because a bare date would misrepresent the
state of the evidence.}
\end{witnesscard}

One orthographic point belongs here rather than in the appendix. The 1962
edition prints \latincue{exspecto resurrectionem mortuorum}; the 1861 hand
missal prints \latincue{expecto}. The word is the same and the sense is the
same; the twentieth-century Vatican editions restored the older Latin spelling
with the internal \latincue{s}. The same restoration accounts for
\latincue{Genetricis} in the Canon and \latincue{neglegentiis} at the offertory,
and it is the single most common way in which the 1962 Latin looks different
from a nineteenth-century printing without being a different text.
